% Компилятор: XeLaTeX. Файл в UTF-8.
% Числа сверены с моделью в коммите 6ee7232.
\documentclass[11pt,a4paper]{article}
\usepackage{fontspec}
\IfFontExistsTF{Arial}{\setmainfont{Arial}}{\setmainfont{Liberation Sans}}
\usepackage{polyglossia}
\setdefaultlanguage{russian}
\usepackage{amsmath,amssymb}
\usepackage[margin=22mm]{geometry}
\usepackage{enumitem}
\usepackage{booktabs,tabularx}
\usepackage[hidelinks]{hyperref}
\setlength{\parindent}{0pt}
\setlength{\parskip}{5pt}
\setlist{itemsep=3pt,topsep=4pt,leftmargin=*}
\setcounter{section}{3}
\newcommand{\rub}{\text{руб.}}

\begin{document}

\section{Где и как проверили решение}

Провели расчётную проверку двух гипотез развития маршрутной сети Red Wings:
нового регионального хаба с пересадками и межрегиональных линий,
объединённых в треугольники из действующих баз в Москве и Екатеринбурге.
Для SJ-100 использовали параметры SSJ-100: 100 кресел и предельную
загрузку 85\% в базовом сценарии.

В лонг-лист вошли 20 городов. После исключения городов ближе 500 км
к Москве или Екатеринбургу в основном рейтинге осталось 15 кандидатов.
Сети рассчитаны для всех 20 городов. В базовом и осторожном сценариях
перебирались лимиты флота
\[
    N_{\max}\in\{4,6,8,10\}\ \text{самолётов}.
\]
Для каждого сценария проверяли варианты без маршрутных субсидий
и с доступной поддержкой по условиям, включённым в модель.
Отдельно оценили чувствительность сети Омска к компенсации стоимости
авиакеросина.

Сегментную модель спроса калибровали по 31 известной паре.
При проверке методом исключения по одному наблюдению (LOOCV)
медианная абсолютная относительная ошибка составила 49,8\% от фактического
потока. Для модели восстановления поездок O\&D ошибка воспроизведения
наблюдаемых сегментных потоков составила 55\%.

Проверили баланс кресел и пассажиров: пассажир с пересадкой занимает
места на двух плечах, выручка его единого билета распределяется между ними.
Спрос общих пар учитывается совместно для всех треугольников сети.
Итоги 84 вариантов и 760 строк маршрутов сверены между CSV, Excel
и интерфейсом.

Проверка выполнена на модели. Коммерческих пилотов и проверки на закрытых
данных продаж авиакомпании не проводили. Стыковочные волны, расписание
и ротации самолётов предстоит проработать отдельно.

\section{На какие данные опираемся}

\subsection*{Данные из источников}

\begin{itemize}
    \item \textbf{Пассажиропотоки.} Сведения о перевозках по парам городов
    за 2025 год: 37 записей в реестре известных потоков, из которых 31
    использована при калибровке; шесть записей использованы для проверки
    нижних границ потока.
    \item \textbf{Транспортная связность.} База проекта «Если быть точным»
    по 180 городам: наличие прямых авиарейсов и железнодорожного сообщения,
    время поездки поездом.
    \item \textbf{Демография и доходы.} Перепись 2021 года, сведения
    Росстата о населении муниципалитетов на 2025 год и зарплатах за 2024 год.
    Население зоны тяготения оценивается в радиусе 75 км от аэропорта.
    \item \textbf{Аэропорты и погода.} Координаты и параметры
    взлётно-посадочных полос из OurAirports; архив погоды NOAA за 2023--2024 годы.
    \item \textbf{Ограничения полётов.} Официальный канал Росавиации:
    5\,981 сообщение за 2026 год в массиве из 8\,329 постов
    за период с 20 октября 2024 года по 9 ноября 2026 года.
    По сообщениям рассчитано число дней с ограничениями «Ковёр».
    \item \textbf{Субсидии и расходы.} Перечень субсидируемых маршрутов
    на 2026 год содержит 446 строк. Для расчёта использованы разбор условий
    программы региональных перевозок, цены авиакеросина, прейскуранты
    аэропортовых сборов и оценки компенсации стоимости керосина Reuters,
    сохранённые в материалах команды.
\end{itemize}

\subsection*{Расчётные оценки и допущения команды}

\textbf{Гравитационная модель спроса.}
Для пар без прямой статистики применена лог-линейная зависимость:
\[
    \ln \widehat Q_{AB}
    =\alpha
    +\beta_1\ln(P_A P_B)
    +\beta_2\ln T_{AB}^{\text{жд}}
    +\gamma_1 I_{AB}^{\text{Москва}}
    +\gamma_2 I_{AB}^{\text{курорт}}.
\]
Здесь $\widehat Q_{AB}$ --- оценка годового потока между городами $A$ и $B$;
$P_A$ и $P_B$ --- численность населения;
$T_{AB}^{\text{жд}}$ --- время поездки поездом;
$I_{AB}^{\text{Москва}}$ и $I_{AB}^{\text{курорт}}$ --- бинарные признаки
столичного и курортного направления. Для сегментного спроса и поездок
O\&D коэффициенты оценены отдельно.

\textbf{Матрица поездок O\&D.}
Поездки распределяются по существующей сети напрямую либо с одной
пересадкой через опорный узел $H$. Для восстановления поездок применяется
ограничение:
\[
    n_{\text{пересадок}}\leq 1,
    \qquad
    \frac{d_{AH}+d_{HB}}{d_{AB}}\leq 1{,}6.
\]
При отборе пересадочных пар для нового хаба предел отношения расстояний
строже: $1{,}35$.

\textbf{Доля захвата спроса.}
На паре без действующего прямого рейса базовая доля зависит от частоты:
\[
    s_{AB}(f)=0{,}65\left(1-e^{-f/3}\right),
    \qquad 3\leq f\leq 14.
\]
Здесь $f$ --- число рейсов в неделю в одну сторону. Для треугольников
верхний предел равен семи. На конкурентных парах доля рассчитывается
по соотношению частот с ограничением сверху; для пересадок используется
отдельная доля захвата.

\textbf{Тарифная функция.}
Базовый тариф билета в одну сторону зависит от расстояния по прямой:
\[
    p(d)=2\,093\,\rub
    +6{,}3\,\frac{\rub}{\text{км}}\cdot d.
\]

\textbf{Параметры самолёта.}
SJ-100 моделируется по аналогии с SSJ-100. При 100 креслах предельная
загрузка составляет 85\% в базе и 80\% в осторожном сценарии.
Расход топлива одного рейса в базовом сценарии:
\[
    F(t)=1{,}55\,\frac{\text{т}}{\text{ч}}\cdot t+0{,}20\,\text{т},
\]
где $t$ --- блок-время рейса в часах. В осторожном сценарии коэффициенты
равны $1{,}90$ т/ч и $0{,}25$ т на рейс.
В основную сеть допускаются плечи до 2\,300 км по прямой,
а для направлений с обходами на юге и западе --- до 2\,150 км.

\textbf{Компенсация стоимости керосина.}
В базовом сценарии суммарная компенсация, включающая демпфер
и обратный акциз, равна 32\,930 руб./т; в осторожном --- 3\,024 руб./т.
Расходы на топливо определяются как
\[
    C_{\text{топл}}=F\left(p_{\text{кер}}-k_{\text{комп}}\right),
\]
где $F$ --- расход топлива, $p_{\text{кер}}$ --- цена керосина,
$k_{\text{комп}}$ --- компенсация на тонну.
Реестр допущений: \path{data/assumptions/network.csv}.

\section{Как решали задачу}

\subsection*{Ход работы}

\begin{enumerate}
    \item \textbf{Подготовили данные.} Составили лонг-лист 20 городов
    и основной рейтинг 15 кандидатов с удалённостью не менее 500 км
    от Москвы и Екатеринбурга. В базовом профиле использованы десять
    критериев с ненулевыми весами. Для оценки спроса собрали данные
    о потоках, населении и связности городов.
    \item \textbf{Рассчитали сети.} Откалибровали модели спроса
    и восстановили матрицу поездок O\&D. Жадным алгоритмом добавляли
    маршруты с наибольшим приростом результата на дополнительный блок-час:
    \[
        j^*=\operatorname*{arg\,max}_{j}
        \frac{\Delta M_j}{\Delta H_j},
    \]
    где $\Delta M_j$ --- прирост выручки за вычетом переменных расходов
    при добавлении маршрута $j$, а $\Delta H_j$ --- дополнительный налёт.
    Проверили лимиты флота 4, 6, 8 и 10 самолётов. Рассчитали билетную
    выручку, субсидии, расходы на топливо, аэропортовое обслуживание,
    экипажи и техническое обслуживание. Такой отбор не гарантирует
    глобально оптимальную сеть.
    \item \textbf{Сравнили варианты.} Сопоставили хабы и прямые
    треугольники при одинаковых ограничениях ресурсов и сценарных
    допущениях. Исключили повторный захват спроса общих рынков и учли
    переход пассажиров с прежних рейсов Red Wings: для базовой сети Омска
    это 8,8 тыс. сегментов в год. Проверили рейтинг в пяти профилях весов
    и чувствительность финансового результата к допущениям.
\end{enumerate}

\subsection*{Связь с учёбой}

\begin{itemize}
    \item \textbf{«Фреймворки Python для анализа данных».}
    Обработка таблиц в pandas и NumPy, оценка коэффициентов спроса
    методом наименьших квадратов, проверка с исключением
    по одному наблюдению (LOOCV).
    \item \textbf{«Алгоритмы и структуры данных».}
    Представление маршрутной сети графом, поиск путей и жадный отбор
    маршрутов при ресурсных ограничениях.
    \item \textbf{«Линейная алгебра».}
    Матричное представление поездок O\&D и их связи с потоками на плечах;
    расчёт ограничений кресельной вместимости.
    \item \textbf{«Технологии промышленного программирования».}
    Модульная архитектура, контроль версий в Git и визуализация
    на JavaScript/D3. Самостоятельно изучили частотные модели выбора
    перевозчика, авиационные тарифы и механизм компенсации керосина.
\end{itemize}

\section{Полученное решение и результат}

\subsection*{Полученное решение}

Создали расчётную модель сети, интерфейс сравнения вариантов
и книгу Excel. Для дальнейшей проработки рекомендуем хаб в Омске:
12 лучей и семь самолётов, включая один резервный.
Омск занимает первое место в четырёх из пяти профилей критериев.
Среднее число дней с ограничениями в аэропортах его сети,
взвешенное по частотам рейсов, составляет 42 дня.
Для сети Сочи этот показатель равен 118 дням.
Для оценки фактического влияния ограничений на пересадки требуется
проверка расписания и длительности закрытий.

\subsection*{Проверка и вывод}

Базовый расчёт выполнен без маршрутных субсидий, при компенсации
керосина 32\,930 руб./т и 52 неделях полётов в году.
Результат указан до аренды самолётов и расходов на создание базы.

\begin{itemize}
    \item \textbf{Хаб Омск, семь самолётов.}
    Всего 577,4 тыс. пассажирских сегментов, из них 568,6 тыс.
    новых для Red Wings; 388,6 тыс. поездок O\&D.
    Выручка с учётом вычета дохода прежних рейсов --- 6,315 млрд руб.,
    результат --- $+503$ млн руб. в год.
    Прирост доли рынка:
    \[
        \Delta s
        =\frac{568\,587}{108\,850\,000}\cdot100
        \approx 0{,}52\ \text{п.п.}
    \]
    \item \textbf{Хаб Сочи, пять самолётов.}
    386,0 тыс. сегментов, выручка 4,337 млрд руб.,
    результат $+569$ млн руб. в год. По финансовому результату
    Сочи выше Омска; показатель ограничений в аэропортах сети также выше.
    \item \textbf{Хаб Уфа, пять самолётов.}
    408,3 тыс. сегментов, выручка 4,139 млрд руб.,
    результат $+252$ млн руб. в год. Уфа исключена из основного
    рейтинга из-за расстояния менее 500 км до Екатеринбурга;
    её сеть рассчитана как дополнительный вариант.
    \item \textbf{Хаб Самара, четыре самолёта.}
    293,0 тыс. сегментов, выручка 2,960 млрд руб.,
    результат $+241$ млн руб. в год.
    \item \textbf{Прямые рейсы, восемь треугольников и восемь самолётов.}
    730,3 тыс. сегментов, выручка 7,221 млрд руб.,
    результат $+44$ млн руб. в год. Положительный результат
    имеют четыре треугольника из восьми.
\end{itemize}

\subsection*{Чувствительность и условия изменения решения}

При полном обнулении компенсации керосина и сохранении базовой сети
результат Омска снижается до $-636$ млн руб. в год.
Если сохраняется обратный акциз 3\,024 руб./т, результат составляет
около $-532$ млн руб. Порог суммарной компенсации для безубыточности
этой сети равен приблизительно 18,4 тыс. руб./т.

В осторожном сценарии без маршрутных субсидий сеть Омска включает
восемь лучей и четыре самолёта, результат составляет
$-2{,}45$ млрд руб. в год. В этом сценарии одновременно используются
нижняя оценка спроса по парам, меньшая доля его захвата, более низкие
тарифы, повышенные расходы и компенсация керосина 3\,024 руб./т.
Годовой спрос распределён на 48 недель полётов.

Рекомендуем поэтапно прорабатывать запуск хаба в Омске:
уточнить спрос по продажам, стоимость базы, условия обслуживания,
доступность самолётов и расписание пересадок. Прибыльные прямые
направления можно рассматривать отдельно с учётом общего парка
и спроса сети.

\subsection*{Файлы расчёта}

Параметры: \path{data/assumptions/network.csv}.
Итоги сетей: \path{data/processed/network_summary.csv}.
Проверка спроса: \path{data/processed/demand_fit.json}
и \path{data/processed/od_fit.json}.
Код: \path{src/network_model.py}.

\end{document}
